\section{微分}\label{sec:ag-3-2}
% 原PDF p44.
\begin{definition}\label{item:ag-3-2-1}
设 $f:X\longrightarrow S$ 是 Noether 概形的分离态射. 此时对角态射 $\Delta_f=\Delta_{X/S}:X\longrightarrow X\times_SX$ 是闭浸入. 称 $X$ 上的凝聚层
\begin{align*}
\Omega_f=\Omega_{X/S}:=\mathcal N_{X/X\times_SX}=\mathcal N_{\Delta_f}
\end{align*}
为 $X$ 在 $S$ 上的\newterm{ag-relative-differentials-sheaf}{相对微分层}{sheaf of relative differentials}，也称\newterm{ag-relative-cotangent-sheaf}{相对余切层}{relative cotangent sheaf}.
\end{definition}
{\small 还应指出，余法层的定义可推广到不一定闭的浸入，余切层的定义则可推广到不一定分离的态射，参见 EGA IV，§16.\par}
% 原PDF p45.
当 $X=\operatorname{Spec}(A)$、$S=\operatorname{Spec}(R)$ 时，对角态射由环 $A$ 的乘法态射 $A\otimes_RA\xrightarrow\mu A$ 给出. 设 $I_{A/R}$ 为此态射的核. 此时 $\Omega_{X/S}$ 由 $A$-模
\begin{align*}
\Omega_{A/R}:=I_{A/R}\otimes_{A\otimes_RA}A=I_{A/R}/I_{A/R}^2
\end{align*}
给出. 注意，$I_{A/R}$ 与 $I_{A/R}^2$ 还自然地以两种方式成为 $A$-模：将 $A$ 映到 $A\otimes1\subset A\otimes_RA$，或将它映到 $1\otimes A$. 在商 $I_{A/R}/I_{A/R}^2$ 上，这两种 $A$-模结构相同，并且与上述模 $I_{A/R}$ 约化所定义的结构一致.
\begin{example}\label{item:ag-3-2-2}
\begin{enumerate}[label=(\roman*)]
\item 设 $A=R[T_1,\ldots,T_n]$. 则 $A\otimes_RA=R[T_1,\ldots,T_n,T'_1,\ldots,T'_n]$，且 $I=(T_1-T'_1,\ldots,T_n-T'_n)$. 容易验证，$\Omega_{A/R}=I/I^2$ 是以 $[(T_i-T'_i)\bmod I^2]$ 为基的自由 $A$-模.
\item 设 $S=\operatorname{Spec}(k)$，其中 $k$ 是代数闭域，$x$ 是 $X$ 的一个闭点. 此时图
\begin{align*}
\begin{tikzcd}[ampersand replacement=\&]
\operatorname{Spec}(k)\arrow[r,hook,"x"]\arrow[d,"x"'] \& X\arrow[d,"{(x,\operatorname{Id})}"]\\
X\arrow[r,hook,"\Delta_{X/k}"'] \& X\times_kX
\end{tikzcd}
\end{align*}
诱导满态射 $\Omega_{X/k,x}\otimes_{\mathcal O_{X,x}}\kappa(x)\twoheadrightarrow\mathfrak m_x/\mathfrak m_x^2$. 它实际上是同构，如下面稍一般的结论所示.\unskip\nobreak\hfill\ensuremath{\diamond}
\end{enumerate}
\renewcommand{\envenddiamond}{}
\end{example}
\begin{proposition}\label{item:ag-3-2-3}
设 $f:X\longrightarrow S$ 是态射，$s:S\hookrightarrow X$ 是 $f$ 的截面. 则典范态射 $s^*\Omega_{X/S}\longrightarrow\mathcal N_{S/X}$ 是同构.
\end{proposition}
\begin{proof}
只须将命题~\ref{item:ag-3-1-5}应用于下图：
\begin{align*}
\begin{tikzcd}[ampersand replacement=\&]
S\arrow[r,hook,"s"]\arrow[d,"s"'] \& X\arrow[r,"f"]\arrow[d,"{(s\circ f,\operatorname{Id})}"] \& S\arrow[d,"s"]\\
X\arrow[r,hook,"\Delta_{X/S}"'] \& X\times_SX\arrow[r,"\pi_1"'] \& X
\end{tikzcd}
\end{align*}
其中 $\pi_1$ 表示第一投影.
\end{proof}
\begin{numpara}[余切层的函子性]\label{item:ag-3-2-4}
注意，当 $f:X\longrightarrow S$ 是浸入时，有 $\Omega_{X/S}=0$，因为此时对角态射是同构. 另一方面，任意交换图
\begin{align}
\begin{tikzcd}[ampersand replacement=\&]
X'\arrow[r,"f'"]\arrow[d,"h"'] \& S'\arrow[d,"g"]\\
X\arrow[r,"f"'] \& S
\end{tikzcd}\tag{3.2.4.1}\label{eq:ag-3-2-4-1}
\end{align}
% 原PDF p46.
均诱导交换图
\begin{align}
\begin{tikzcd}[ampersand replacement=\&]
X'\arrow[r,hook,"\Delta_{f'}"]\arrow[d,"h"'] \& X'\times_{S'}X'\arrow[d,"h\times h"]\\
X\arrow[r,hook,"\Delta_f"'] \& X\times_SX
\end{tikzcd}\tag{3.2.4.2}\label{eq:ag-3-2-4-2}
\end{align}
从而给出一个~\eqref{eq:ag-3-1-4-2}型的典范态射
\begin{align}
h^*\Omega_f\longrightarrow\Omega_{f'}.\tag{3.2.4.3}\label{eq:ag-3-2-4-3}
\end{align}
当 $h$ 是开浸入且 $g$ 是 $S$ 的恒等态射时，它是同构. 事实上，此时图~\eqref{eq:ag-3-2-4-2}是笛卡尔图，可以应用分目~\ref{item:ag-3-1-4}. 这特别说明，$\Omega_{X/S}$ 与局部化相容，因而只须在仿射坐标图中计算它. 更确切地说，若 $U=\operatorname{Spec}(A)$ 是 $X$ 的仿射开集，且 $f|_U$ 经过 $S$ 的仿射开集 $V=\operatorname{Spec}(R)$ 分解，则 $\Omega_{X/S}(U)=\Omega_{A/R}$. 若 $x$ 是 $X$ 的一点，其在 $S$ 中的像为 $s$，则 $\Omega_{X/S,x}=\Omega_{\mathcal O_{X,x}/\mathcal O_{S,s}}$.
下面是上述态射为同构的另一种情形.
\end{numpara}
\begin{proposition}\label{item:ag-3-2-5}
若图~\eqref{eq:ag-3-2-4-1}是笛卡尔图，则态射~\eqref{eq:ag-3-2-4-3}是同构.
\end{proposition}
\begin{proof}
只须将图~\eqref{eq:ag-3-2-4-2}补全如下：
\begin{align*}
\begin{tikzcd}[ampersand replacement=\&]
X'\arrow[r,hook,"\Delta_{f'}"]\arrow[d,"h"'] \& X'\times_{S'}X'\arrow[r,"\pi_1"]\arrow[d,"h\times h"] \& X'\arrow[d,"h"]\\
X\arrow[r,hook,"\Delta_f"'] \& X\times_SX\arrow[r,"\pi_1"'] \& X
\end{tikzcd}
\end{align*}
并应用命题~\ref{item:ag-3-1-5}的 (ii).
\end{proof}
为说明与通常微分直观的联系，先回顾导子的概念.
\begin{definition}\label{item:ag-3-2-6}
设 $R$ 是交换环，$A$ 是交换 $R$-代数. 取值于 $A$-模 $M$ 的 $R$-\newterm{ag-derivation}{导子}{derivation} $d:A\longrightarrow M$ 是满足下列条件的 $R$-线性映射：
\begin{itemize}
\item 对任意 $a,a'\in A$，有 $d(aa')=ad(a')+d(a)a'$.
\item 对任意 $r\in R$，有 $d(r)=0$.
\end{itemize}
\end{definition}
$A$ 到 $M$ 的所有 $R$-导子构成的集合 $\operatorname{Der}_R(A,M)$ 自然是 $A$-模. 用公式
\begin{align*}
d_{A/R}(a):=(a\otimes1-1\otimes a)\bmod I_{A/R}^2
\end{align*}
定义一个取值于 $A$-模 $\Omega_{A/R}$ 的 $R$-导子.
% 原PDF p47.
二元组 $(\Omega_{A/R},d_{A/R})$ 在配备 $R$-导子的 $A$-模中具有如下意义下的万有性.
\begin{proposition}\label{item:ag-3-2-7}
二元组 $(\Omega_{A/R},d_{A/R})$ 表示范畴 $\operatorname{Mod}(A)$ 上的函子 $M\mapsto\operatorname{Der}_R(A,M)$. 换言之，与 $d_{A/R}$ 复合，对任意 $A$-模 $M$ 给出同构 $\operatorname{Hom}_A(\Omega_{A/R},M)\xrightarrow\sim\operatorname{Der}_R(A,M)$.
\end{proposition}
\begin{proof}
设 $d:A\longrightarrow M$ 是 $R$-导子. 记 ${}_A(A\otimes_RA)$ 为 $A$ 通过左作用赋予 $A\otimes_RA$ 的 $A$-模结构. $R$-线性映射 $d$ 诱导 $A$-线性映射 ${}_Ad:{}_A(A\otimes_RA)\longrightarrow M$. 利用 $d$ 的导子性质计算可知，${}_Ad$ 在 $A$-子模 $I_{A/R}^2$ 上为零. 因而存在 $A$-模态射 $f:\Omega_{A/R}\longrightarrow M$，使图
\begin{align*}
\begin{tikzcd}[ampersand replacement=\&]
{}_A(A\otimes_RA)\arrow[r,"{}_Ad"] \& M\\
I_{A/R}\arrow[u,hook]\arrow[r,two heads] \& \Omega_{A/R}\arrow[u,dashed,"f"']
\end{tikzcd}
\end{align*}
交换. 当 $(M,d)=(\Omega_{A/R},d_{A/R})$ 时，有 $f=\operatorname{Id}_{\Omega_{A/R}}$，命题由此得出.
\end{proof}
\begin{example}\label{item:ag-3-2-8}
\begin{enumerate}[label=(\roman*)]
\item 设 $A$ 是代数闭域上的有限生成代数，$x$ 是 $\operatorname{Spec}(A)$ 的闭点，其理想为 $\mathfrak m_x$. 映射
\begin{align*}
A&\longrightarrow\mathfrak m_x/\mathfrak m_x^2,\\
f&\longmapsto d_x(f):=[(f-f(x))\bmod\mathfrak m_x^2]
\end{align*}
是一个 $k$-导子，称为\emph{在 $x$ 处的微分}. 所诱导的态射 $\Omega_{A/k}\otimes_A(A/\mathfrak m_x)\longrightarrow\mathfrak m_x/\mathfrak m_x^2$ 正是例~\ref{item:ag-3-2-2}中的态射. 特别地，当 $A=k[T_1,\ldots,T_n]$ 时，$d_x(T_i)$ 构成 $\mathfrak m_x/\mathfrak m_x^2$ 的一组基，并且有 $d_x(f)=\sum_i\frac{\partial f}{\partial T_i}(x)d_xT_i$，如例~\eqref{eq:ag-3-1-13-1}所示.
\item 当 $A=R[T_1,\ldots,T_n]$ 时，例~\ref{item:ag-3-2-2}说明，$\Omega_{A/R}$ 是以 $d_{A/R}(T_i)$ 为基的自由模. 此时导子 $d_{A/R}$ 由下式给出：
\begin{align*}
d_{A/R}:\quad A&\longrightarrow\Omega_{A/R},\\
f&\longmapsto d_{A/R}(f):=\sum_i\frac{\partial f}{\partial T_i}d_{A/R}(T_i).
\tag*{$\diamond$}
\end{align*}
\end{enumerate}
\renewcommand{\envenddiamond}{}
\end{example}
\begin{proposition}\label{item:ag-3-2-9}
设 $Y\xrightarrow gX\xrightarrow fS$ 是（Noether）概形的两个（分离）态射. 则在 $Y$ 上存在凝聚 $\mathcal O_Y$-模的正合列：
\begin{align*}
g^*\Omega_f\longrightarrow\Omega_{f\circ g}\longrightarrow\Omega_g\longrightarrow0,
\end{align*}
其中所有箭头均为~\eqref{eq:ag-3-2-4-3}型. 此外，若对角浸入 $\Delta_g$ 正则（此时称 $g$ \newterm{ag-differentially-smooth}{微分光滑}{differentially smooth}），且 $g$ 平坦，则在左端添加零态射后，此列在左端也正合.
\end{proposition}
% 原PDF p48.
\begin{proof}
将命题~\ref{item:ag-3-1-15}应用于浸入 $Y\hookrightarrow Y\times_XY\hookrightarrow Y\times_SY$，得到正合列
\begin{align}
\Delta_{Y/X}^*(\mathcal N_{Y\times_XY/Y\times_SY})\longrightarrow\Omega_{Y/S}\longrightarrow\Omega_{Y/X}\longrightarrow0.\tag{3.2.9.1}\label{eq:ag-3-2-9-1}
\end{align}
另一方面，交换图
\begin{align*}
\begin{tikzcd}[ampersand replacement=\&]
Y\times_XY\arrow[r,hook]\arrow[d] \& Y\times_SY\arrow[d]\\
X\arrow[r,hook] \& X\times_SX
\end{tikzcd}
\end{align*}
是笛卡尔图，从而诱导满态射 $g^*\Omega_{X/S}\longrightarrow\Delta_{Y/X}^*(\mathcal N_{Y\times_XY/Y\times_SY})$. 命题中的列由此正合. 若还假设 $g$ 平坦，则上述满态射是双射. 再假设 $\Delta_{Y/X}$ 正则，则仍由命题~\ref{item:ag-3-1-15}，列~\eqref{eq:ag-3-2-9-1}在左端正合.
\end{proof}
\begin{remark}\label{item:ag-3-2-10}
当 $Y=\mathbb A^n\times X=\operatorname{Spec}(\mathcal O_X[T_1,\ldots,T_n])$（$X$ 上的相对谱）时，或者当 $Y$ 是 $X$ 上的向量丛时，假设``$g$ 平坦且微分光滑''特别地成立. 我们将看到，这是光滑态射的特征性质之一.\unskip\nobreak\hfill\ensuremath{\diamond}
\renewcommand{\envenddiamond}{}
\end{remark}
\begin{example}\label{item:ag-3-2-11}
考虑 $Y=\mathbb A_k^n$、$X=\mathbb A_k^m$、$S=\operatorname{Spec}(k)$ 的情形. 命题~\ref{item:ag-3-2-9}的正合列使我们可将 $\Omega_{Y/X}$ 计算为 $g^*\Omega_{X/k}\longrightarrow\Omega_{Y/k}$ 的余核. 来具体写出此态射. 态射 $g$ 由环态射 $k[X_1,\ldots,X_m]\longrightarrow k[T_1,\ldots,T_n]$ 给出；后者由它在 $X_j$ 上的取值确定，这些取值是 $T_i$ 的多项式 $P_j$. 由例~\ref{item:ag-3-2-2}，$g^*\Omega_{X/k}$ 是以 $d_{X/k}(X_j)$（$j=1,\ldots,m$）为基的自由 $k[T_1,\ldots,T_n]$-模，而 $\Omega_{Y/k}$ 是以 $d_{Y/k}(T_i)$（$i=1,\ldots,n$）为基的自由 $k[T_1,\ldots,T_n]$-模. 所求映射于是由下式给出：
\begin{align*}
g^*\Omega_{X/k}&\longrightarrow\Omega_{Y/k},\\
d_{X/k}(X_j)&\longmapsto\sum_{i=1}^n\frac{\partial P_j}{\partial T_i}d_{Y/k}(T_i).
\tag*{$\diamond$}
\end{align*}
\renewcommand{\envenddiamond}{}
\end{example}
当 $g$ 是闭浸入时，命题~\ref{item:ag-3-2-9}中正合列的右项为零. 此时可向左延长该正合列.
\begin{proposition}\label{item:ag-3-2-12}
设 $Y\xrightarrow gX\xrightarrow fS$ 是（Noether）概形的两个（分离）态射，其中 $g$ 是闭浸入，其理想为 $\mathcal I\subset\mathcal O_X$. 则在 $Y$ 上有正合列
\begin{align*}
\mathcal N_g\xrightarrow\delta g^*\Omega_f\longrightarrow\Omega_{fg}\longrightarrow0,
\end{align*}
其中 $\delta$ 由将导子 $d_{X/S}:\mathcal O_X\longrightarrow\Omega_f$ 限制到 $\mathcal I$，再模 $\mathcal I$ 约化而定义.
\end{proposition}
% 原PDF p49.
\begin{proof}
在 $\Omega_{fg}$ 处的正合性来自前一命题以及 $\Omega_g=0$ 这一事实. 对于其余部分，归结到仿射情形：$S=\operatorname{Spec}(R)$，$X=\operatorname{Spec}(A)$，$Y=\operatorname{Spec}(A/I)$. 由命题~\ref{item:ag-3-2-7}，只须证明，对任意 $A/I$-模 $M$，列
\begin{align*}
0\longrightarrow\operatorname{Der}_R(A/I,M)\longrightarrow\operatorname{Der}_R(A,M)\longrightarrow\operatorname{Hom}_A(I,M)
\end{align*}
在中间正合，这是显然的.
\end{proof}
\begin{example}\label{item:ag-3-2-13}
\begin{enumerate}[label=(\roman*)]
\item 取 $S=\operatorname{Spec}(k)$，其中 $k$ 是代数闭域，并取 $Y=\{x\}$. 此时 $\Omega_{Y/S}=0$，从而有满态射 $\mathfrak m_x/\mathfrak m_x^2\xrightarrow\delta\Omega_{X/k}\otimes\kappa(x)$. 它是同构，与例~\ref{item:ag-3-2-2}(ii) 和例~\ref{item:ag-3-2-8}(i) 中的同构互逆.
\item 仍设 $S=\operatorname{Spec}(k)$，并进一步假设 $Y$ 是 $X=\mathbb A_k^n$ 的子簇，其代数为
\begin{align*}
k[Y]:=k[T_1,\ldots,T_n]/(P_1,\ldots,P_r).
\end{align*}
命题~\ref{item:ag-3-2-12}的正合列以及 $\Omega_{X/k}=k[X]^n$ 这一已知结果，使我们可将 $\Omega_{Y/k}$ 计算为映射 $k[Y]^r\longrightarrow k[Y]^n$ 的余核，此映射的矩阵为
\begin{align*}
\biggl(\left.\frac{\partial P_j}{\partial T_i}\right|_Y\biggr)_{i=1,\ldots,n;\,j=1,\ldots,r}.
\end{align*}
\item 设 $k$ 是域，$\kappa=k[T]/P(T)$ 是 $k$ 的有限单生成域扩张. 因而，$\kappa$ 是 $\mathbb A_k^1$ 中某个闭点 $x$ 的剩余域. $\kappa$-向量空间正合列
\begin{align*}
\mathfrak m_x/\mathfrak m_x^2\longrightarrow\Omega_{\mathbb A^1/k,x}\longrightarrow\Omega_{\kappa/k}\longrightarrow0
\end{align*}
同构于 $\kappa\xrightarrow{\times P'(x)}\kappa\longrightarrow\Omega_{\kappa/k}$，由此可见，$\Omega_{\kappa/k}$ 为零当且仅当 $P'(x)\ne0$，即当且仅当 $\kappa$ 在 $k$ 上可分.\unskip\nobreak\hfill\ensuremath{\diamond}
\end{enumerate}
\renewcommand{\envenddiamond}{}
\end{example}
\begin{remark}\label{item:ag-3-2-14}
命题~\ref{item:ag-3-2-9}与~\ref{item:ag-3-2-12}中的部分正合列是一个更一般现象的表现. 事实上，当 $X$ 在 $S$ 上不光滑时，余切层 $\Omega_{X/S}$ 不是足够精细的不变量. 用一个``\newterm{ag-cotangent-complex}{余切复形}{cotangent complex}''替代它，上述部分正合列便是从余切复形的``正合三角''所伴随的长正合列中截取的（参见 Illusie）.\unskip\nobreak\hfill\ensuremath{\diamond}
\renewcommand{\envenddiamond}{}
\end{remark}
\begin{example}\label{item:ag-3-2-15}
在命题~\ref{item:ag-3-2-12}的情形下，设 $X=\operatorname{Spec}(\mathcal O_S[T_1,\ldots,T_n])$，$Y=\operatorname{Spec}(\mathcal O_S[T_1,\ldots,T_m])$. 则 $Y$ 在 $X$ 中的理想 $\mathcal I$ 是由 $T_{m+1},\ldots,T_n$ 生成的理想. 此时命题~\ref{item:ag-3-2-12}中的列在左端正合，并同构于列
\begin{align*}
0\longrightarrow\bigoplus_{i=m+1}^n\mathcal O_Y\overline T_i
\xrightarrow\delta\bigoplus_{i=1}^n\mathcal O_Yd_{X/S}T_i
\longrightarrow\bigoplus_{i=1}^m\mathcal O_Yd_{Y/S}T_i\longrightarrow0,
\end{align*}
其中 $\overline T_i$ 是 $T_i\in\mathcal I$ 在 $\mathcal I/\mathcal I^2=\mathcal N_{Y/X}$ 中的像，$\delta$ 将 $\overline T_i$ 映到 $d_{X/S}(T_i)$.\unskip\nobreak\hfill\ensuremath{\diamond}
\renewcommand{\envenddiamond}{}
\end{example}
\begin{remark}\label{item:ag-3-2-16}
在例~\ref{item:ag-3-1-13}的情形下，$Y\subset X$ 是域 $k$ 上簇的一个光滑对，命题~\ref{item:ag-3-2-12}中的列在左端也正合，且其中各项在 $Y$ 上均局部自由. 事实上，已经看到浸入 $Y\hookrightarrow X$ 正则，所以 $\mathcal N_{Y/X}$ 局部自由，秩为 $\operatorname{codim}(Y,X)$. 另一方面，$Y\times_kY$ 也是光滑簇，因此对角浸入 $\Delta_{Y/k}$ 正则，而 $\Omega_{Y/k}$ 局部自由，秩为 $\dim(Y)$. 最后，$i^*\Omega_{X/k}$ 局部自由，秩为 $\dim(X)$. 秩的可加性迫使左端正合. 这样便确认了余法层在正则点处的几何意义：
% 原PDF p50上部.
它是子簇的切空间在环境余切空间中的正交补. 在``光滑态射''一节中，将在任意基概形 $S$ 上推广这一结论.\unskip\nobreak\hfill\ensuremath{\diamond}
\renewcommand{\envenddiamond}{}
\end{remark}
